\section{Evaluation}
\label{sec:eval}

Four research questions organize the evaluation. \textbf{RQ1}: what does making
the graph admissible cost in task accuracy? \textbf{RQ2}: is the retrofit
actually admissible under fallback-free compilation, and what is the latency and
throughput price at the engine boundary? \textbf{RQ3}: how does that price
appear at the application boundary once preprocessing, layout conversion, and
host decoding are accounted for? \textbf{RQ4}: what does strict placement return
in energy per frame and in GPU availability under contention? RQ1--RQ3 are
answered with measurements; RQ4 is specified, instrumented, and not yet
measured, and we mark it as such rather than argue it.

We test one hypothesis, H1: a substitution restricted to the inadmissible
regions of a strong detector, initialized from that detector's weights, retains
most of its accuracy, so that admissibility can be bought at a task cost small
enough to be irrelevant to the deployment decision. The deployment decision
itself is then determined by RQ2--RQ4 and not by H1.

\subsection{Setup}
\label{sec:setup}

\paragraph*{Platform}
All device measurements come from one NVIDIA Jetson Orin module. The build logs
record device name \texttt{Orin}, compute capability $8.7$, 16 streaming
multiprocessors, $62\,840$~MiB of shared device memory, a 256-bit memory bus, an
application compute clock of $1.3$~GHz, TensorRT $10.3.0$, and DLA core~0 with
DLA memory pools of 1~MiB managed SRAM, 1024~MiB local DRAM, and 512~MiB global
DRAM. \TODO{exact module SKU and carrier, JetPack/L4T version, CUDA and cuDNN
versions, active \texttt{nvpmodel} power mode, \texttt{jetson\_clocks} state,
measured GPU/DLA/EMC clocks during each run, and thermal state; none of these
are in the current logs and all must be recorded with the final runs}

\paragraph*{Models and accuracy protocol}
The reference is the official COCO-pretrained YOLO11n checkpoint; the retrofit
is its best validation checkpoint \TODO{exact training schedule of the reported
retrofit checkpoint: epoch budget of each phase and the frozen-adaptation
length; the checkpoint file name in our records suggests a different epoch
count than the one we had noted}. Both are evaluated on the same COCO val2017
split \cite{lin2014coco} at $640\times640$ over the same 80 classes. Precision
and recall are averaged over classes at the confidence operating point that
maximizes mean F1. Reference retention is $100\,Q(G)/Q(G_0)$ with
$Q=\mathrm{mAP}_{50\text{--}95}$.

\paragraph*{Deployment protocol}
Three builds are used: TensorRT GPU FP16 as the primary baseline, TensorRT with
DLA requested and GPU fallback enabled as a partitioning diagnostic, and strict
DLA FP16 with fallback disabled. All comparisons are at FP16, batch one, and
$640\times640$. Engine-boundary numbers come from \texttt{trtexec} with
generated inputs, a 2~s warm-up, a 30~s run, one inference stream, spin waiting,
and H2D and D2H transfers included; TensorRT's reported host latency is the sum
of H2D, device execution, and D2H, and excludes preprocessing and detector
post-processing, so we call it engine-plus-transfer latency and never
end-to-end latency. Application-level numbers come from one instrumented tester
that letterboxes a $960\times1280$ image to $640\times640$ and runs 100 timed
iterations after ten full-pipeline warm-ups, with single-threaded CPU
post-processing, $\tau=0.25$, and NMS IoU $0.45$; file reading, annotation, and
image writing are excluded.

\subsection{RQ1: Accuracy and Retention (the Price of Admission)}
\label{sec:rq1}

Table~\ref{tab:accuracy} reports the model-level comparison. The retrofit
reaches $37.68\%$ $\mathrm{mAP}_{50\text{--}95}$ against $39.40\%$ for the
reference, i.e.\ it gives up $1.72$ points and retains $95.6\%$; precision,
recall, and $\mathrm{mAP}_{50}$ fall by $0.94$, $1.22$, and $1.97$ points. This
supports H1 and closes RQ1: the accuracy cost of removing two inadmissible
subgraphs from a modern detector is under two points of
$\mathrm{mAP}_{50\text{--}95}$, which for most robot deployments is not the
term that decides where the model runs. We deliberately do not present this as
the result of the paper. It is the precondition that makes the rest of the
comparison meaningful, and it is measured on the PyTorch path, not through the
serialized accelerator engine.

\begin{table}[t]
\caption{COCO val2017 detection accuracy at $640\times640$, percent. The first
four columns are measured on the ordinary PyTorch inference path. The last
column, which is what the deployed strict engine actually delivers, is not yet
measured.}
\label{tab:accuracy}
\centering
\scriptsize
\setlength{\tabcolsep}{2.2pt}
\begin{tabular}{lccccc}
\toprule
 & \multicolumn{4}{c}{PyTorch path} & Strict DLA engine \\
\cmidrule(lr){2-5}\cmidrule(lr){6-6}
Model & P & R & $\mathrm{mAP}_{50}$ & $\mathrm{mAP}_{50\text{--}95}$
      & $\mathrm{mAP}_{50\text{--}95}$ \\
\midrule
YOLO11n      & 65.64 & 50.25 & 55.14 & 39.40 & n/a (build rejected) \\
YOLO11-DLA-n & 64.70 & 49.03 & 53.17 & 37.68 & \TODO{measure} \\
\bottomrule
\end{tabular}
\end{table}

The empty cell in Table~\ref{tab:accuracy} is a real gap, not a formality.
\TODO{COCO $\mathrm{mAP}_{50\text{--}95}$ evaluated from the serialized strict
DLA engine. The stock Ultralytics TensorRT backend cannot bind CHW16 outputs or
a DLA image input, so this requires running the validator through the custom
tester path of Sec.~\ref{sec:packing}; FP16 rounding, the compiler's DLA
softmax approximation, and the packing path can all move the number} Until it is
filled, every accuracy statement in this paper describes the PyTorch model and
not the engine.

\subsection{RQ2: Admissibility and the Engine Boundary}
\label{sec:rq2}

\paragraph*{Admissibility}
The unmodified detector is rejected. With DLA as the default device and fallback
disabled, the build aborts at \texttt{/model.10/m/m.0/attn/MatMul} with
\emph{``MatMul is unsupported on DLA''} followed by a DLA validation failure, so
no engine exists to measure. The retrofit is accepted: the verbose build reports
299 layers running on DLA and \emph{``No layer is running on GPU''}, and the
engine inspector represents the whole learned graph as a single \texttt{DLA}
layer with the three declared prediction bindings. Compliance gates 1--4 of
Sec.~\ref{sec:deployment} are therefore satisfied; gate~5 is satisfied at the
level of the interface check reported in Sec.~\ref{sec:rq3} and not yet at the
level of dataset accuracy.

\paragraph*{Why partial offloading is not an option}
The fallback build is the natural compromise a practitioner would try, and it
does not work. TensorRT keeps 317 layers on DLA and moves 17 to the GPU: the two
attention matrix multiplications, eight reshape/transpose shuffles, two slices,
one division, and four constants. The reasons the log gives for the head are
instructive, because they are shape reasons rather than operator reasons:
\emph{``Input tensor has less than 4 dimensions''} and \emph{``Dimension: 3
(8400) exceeds maximum allowed size for DLA: 8192''}. The decoded interface is
inadmissible twice over, once for rank and once for extent, which is what
motivates the boundary of Sec.~\ref{sec:boundary}.

After fusion the resulting engine contains 47 nodes: four DLA loadables, 11 GPU
compute kernels, 22 reformat nodes, seven no-ops, and three constants. The layer
profiler attributes $93.97\%$ of execution time to the four DLA loadables,
$4.73\%$ to the reformat nodes, and $1.31\%$ to all GPU compute together; the
largest single DLA segment, which ends at the attention transpose, accounts for
$42.54\%$ ($14.581$~ms). So the GPU does almost no arithmetic in this build, and
the build is still the slowest of the three at $34.678$~ms median. Coverage
statistics are the wrong metric for a partitioned graph: what costs time is the
partitioning itself, in layout conversion and in serialization of four
loadables. A build that is $94\%$ on the accelerator still holds the GPU, still
pays $22.80$~MiB of execution-context device memory, and its median
engine-plus-transfer latency is $26.0\%$ above that of the fully placed strict
engine. We report it as a diagnostic and not as a
deployment option.

\paragraph*{The price at the engine boundary}
Table~\ref{tab:engine} gives the matrix. The comparison that matters is the
strict engine against the primary baseline, the reference detector on the GPU:
$27.515$~ms against $3.740$~ms median engine-plus-transfer latency, a factor of
$7.36$, and $37.2$ against $288.3$ queries/s saturated throughput, i.e.\
$12.9\%$. Against the retrofit's own GPU engine the factor is $7.86$. Transfers
are not the reason: for the strict engine the median splits into $0.393$~ms
H2D, $26.816$~ms device execution, and $0.305$~ms D2H, so $97.5\%$ of the
engine-boundary time is the accelerator computing. This is what a
fixed-function core costs on this workload, and no boundary placement or
implementation effort inside the engine will remove it.

Two secondary observations. First, the retrofit is not slower than the reference
on the GPU: $3.500$ against $3.740$~ms median and $317.2$ against
$288.3$~queries/s, which is an interface difference rather than an architecture
win, since the retrofit's engine stops before the decode that the reference
performs internally. Second, the strict engine is the cheapest to build and to
hold: $18.7$~s build time against $471.7$~s for the reference GPU engine,
$6.52$~MiB serialized against $8.15$~MiB, and zero execution-context device
memory against $9.47$~MiB, because its activations live in the accelerator's own
memory pools. Those allocator figures support the placement audit; they are not
a substitute for the concurrency measurement of Sec.~\ref{sec:rq4}.

\begin{table}[t]
\caption{FP16 TensorRT engine boundary, batch one, $640\times640$, DLA core~0.
Latency includes H2D, device execution, and D2H, and excludes pre- and
post-processing. DLA share is profiler time, not layer count. Each row is one
30~s \texttt{trtexec} launch.}
\label{tab:engine}
\centering
\scriptsize
\setlength{\tabcolsep}{2.0pt}
\begin{tabular}{lccrrr}
\toprule
Model / execution & Build & DLA & Median & P95 & Thr. \\
                  &       & (\%) & (ms)  & (ms) & (q/s) \\
\midrule
YOLO11n / GPU (baseline)   & yes & 0    & 3.740  & 3.753  & 288.3 \\
YOLO11n / DLA+GPU (diagn.) & yes & 94.0 & 34.678 & 34.942 & 29.2 \\
YOLO11n / strict DLA       & \textbf{no} & --- & --- & --- & --- \\
YOLO11-DLA-n / GPU         & yes & 0    & 3.500  & 3.512  & 317.2 \\
YOLO11-DLA-n / strict DLA  & yes & 100  & 27.515 & 27.718 & 37.2 \\
\bottomrule
\end{tabular}
\end{table}

\subsection{RQ3: Application-Level Accounting}
\label{sec:rq3}

% Stacked per-stage accounting; bar lengths are the median stage times of
% Table~\ref{tab:e2e}. Horizontal unit: 1 ms.
\begin{figure}[t]
\centering
\begin{tikzpicture}[x=0.090cm, y=1cm,
  lab/.style={font=\scriptsize, anchor=east},
  tot/.style={font=\scriptsize, anchor=west}]

\definecolor{cpre}{RGB}{190,190,190}
\definecolor{cio}{RGB}{235,180,60}
\definecolor{ceng}{RGB}{110,145,205}
\definecolor{cpost}{RGB}{140,180,120}

\newcommand{\stackbar}[6]{%
  \node[lab] at (-1,#1+0.21) {#6};
  \fill[cpre]  (0,#1) rectangle (#2,#1+0.42);
  \fill[cio]   (#2,#1) rectangle (#2+#3,#1+0.42);
  \fill[ceng]  (#2+#3,#1) rectangle (#2+#3+#4,#1+0.42);
  \fill[cpost] (#2+#3+#4,#1) rectangle (#2+#3+#4+#5,#1+0.42);
  \draw[very thin] (0,#1) rectangle (#2+#3+#4+#5,#1+0.42);
}

\stackbar{0}{4.900}{2.843}{8.153}{3.250}{11n / GPU}
\node[tot] at (19.15,0.21) {19.203};
\stackbar{-0.60}{4.572}{3.848}{7.560}{3.717}{11-DLA-n / GPU}
\node[tot] at (19.70,-0.39) {19.831};
\stackbar{-1.20}{8.617}{4.002}{35.198}{4.866}{11n / DLA+GPU}
\node[tot] at (52.69,-0.99) {52.782};
\stackbar{-1.80}{8.946}{13.992}{27.788}{5.711}{11-DLA-n / DLA}
\node[tot] at (56.44,-1.59) {56.497};

\draw[->] (0,-2.05) -- (64,-2.05);
\foreach \x in {0,10,20,30,40,50,60}
  {\draw (\x,-2.00) -- (\x,-2.12) node[below, font=\scriptsize] {\x};}
\node[font=\scriptsize, anchor=north] at (32,-2.42) {median stage latency (ms)};

\fill[cpre] (0,-2.95) rectangle (3,-2.79);
\node[font=\scriptsize, anchor=west] at (3.6,-2.87) {pre};
\fill[cio] (12,-2.95) rectangle (15,-2.79);
\node[font=\scriptsize, anchor=west] at (15.6,-2.87) {I/O};
\fill[ceng] (24,-2.95) rectangle (27,-2.79);
\node[font=\scriptsize, anchor=west] at (27.6,-2.87) {engine};
\fill[cpost] (42,-2.95) rectangle (45,-2.79);
\node[font=\scriptsize, anchor=west] at (45.6,-2.87) {post};
\end{tikzpicture}
\caption{Per-stage accounting of the repeated-image pipeline check, reference
implementation, single-threaded CPU pre- and post-processing
(Table~\ref{tab:e2e}). I/O is input packing plus H2D, plus D2H plus output
unpacking. Numbers at the right are median totals; because stage medians are
computed independently they need not sum exactly to the median total. The
strict DLA row is the only one whose I/O stage carries a DLA-native layout
conversion, and the only one whose engine stage uses no GPU.}
\label{fig:stack}
\end{figure}


Table~\ref{tab:e2e} and Fig.~\ref{fig:stack} decompose one instrumented
pipeline over the four configurations that have data. The strict path measures
$56.497$~ms median total against $19.203$~ms for the baseline, a factor of
$2.94$; the sequential image rates are $17.70$ and $52.08$~images/s, which are
inverse median latencies and not saturated throughput. Decomposing the
$37.29$~ms difference by stage is more informative than the ratio: the engine
contributes $+19.64$~ms, layout conversion $+11.15$~ms, preprocessing
$+4.05$~ms, and host decoding $+2.46$~ms.

Only the first of those four terms is intrinsic to the accelerator. The second
is the reference implementation's cost, and we state its nature plainly:
packing the DLA image input and unpacking the three CHW16 outputs together take
$13.992$~ms ($7.755$~ms in, $6.237$~ms out) against $2.843$~ms for the
baseline's linear bindings, for buffers of $3.28$ and $2.42$~MB respectively.
That is an effective conversion rate of roughly $0.42$ and $0.39$~GB/s, far
below the module's memory bandwidth, which is unsurprising for
single-threaded PyTorch tensor permutations on the CPU as described in
Sec.~\ref{sec:packing}. Table~\ref{tab:e2e} is therefore a characterization of
\emph{this reference implementation}, not a measurement of the intrinsic cost of
the DLA path, and the comparison it supports is bounded pending an optimized
converter. \TODO{stage row for an optimized I/O path: a CUDA or NEON packer, a
pinned staging buffer, conversion overlapped with execution, and a decoder that
consumes CHW16 directly; this row is the honest upper bound on how much of the
$11.15$~ms is removable}

The third and fourth terms are not explained by any code path. Preprocessing is
identical OpenCV letterboxing of the identical image in every row, yet it
measures $4.900$ and $4.572$~ms in the two GPU rows and $8.617$ and $8.946$~ms
in the two DLA rows. Host post-processing is likewise the same single-threaded
code on the same 13 surviving anchors for both retrofit rows, yet it measures
$3.717$~ms with the GPU engine and $5.711$~ms with the strict engine. The two
DLA-row processes are thus about $6.4$~ms slower on stages that never touch the
accelerator. We report this as an unexplained cross-process measurement
difference between process launches and decline to attribute it to thermal,
clock, or scheduling effects, none of which we recorded. It bounds the precision
of the totals in Table~\ref{tab:e2e} and is a first-order reason to repeat the
experiment as specified in Sec.~\ref{sec:validity}.

Two further results belong here. The interface check passes: the retrofit's GPU
and strict-DLA executions return the same two detections with the same
confidences (\texttt{person} $0.871$, \texttt{bench} $0.317$) and box
coordinates differing by at most $0.3$~pixels, which validates the packing and
unpacking path for this input without substituting for engine-level dataset
validation. And the confidence-first decoder of Sec.~\ref{sec:deployment} is
what keeps host decoding affordable: it reduces median post-processing from
$17.011$ to $3.717$~ms on the GPU engine and from $26.744$ to $5.711$~ms on the
strict engine, while returning the same detections. Without it, the host decode
alone would exceed the entire baseline pipeline.

\begin{table}[t]
\caption{Repeated-image pipeline check, one process launch per row, medians in
milliseconds; I/O sums input packing/H2D and D2H/unpacking; P95 and P99 refer to
the total. \emph{This is a characterization of the reference implementation}:
pre-, post-processing and layout conversion are single-threaded CPU code that
has not been optimized.}
\label{tab:e2e}
\centering
\scriptsize
\setlength{\tabcolsep}{1.6pt}
\begin{tabular}{lrrrrrrr}
\toprule
Execution & Pre & I/O & Eng. & Post & Total & P95 & P99 \\
\midrule
11n / GPU        & 4.900 & 2.843  & 8.153  & 3.250 & \textbf{19.203} & 19.574 & 19.880 \\
11n / DLA+GPU    & 8.617 & 4.002  & 35.198 & 4.866 & \textbf{52.782} & 53.098 & 53.396 \\
11-DLA-n / GPU   & 4.572 & 3.848  & 7.560  & 3.717 & \textbf{19.831} & 20.343 & 20.544 \\
11-DLA-n / DLA   & 8.946 & 13.992 & 27.788 & 5.711 & \textbf{56.497} & 56.864 & 57.781 \\
\addlinespace
\multicolumn{8}{l}{\emph{optimized I/O bound, same engine:}} \\
11-DLA-n / DLA   & \multicolumn{7}{l}{\TODO{optimized packer}} \\
\bottomrule
\end{tabular}
\end{table}

\subsection{RQ4: Energy and GPU Isolation}
\label{sec:rq4}

\begin{figure}[t]
\centering
\fbox{\parbox[c][3.1cm][c]{0.94\columnwidth}{\centering\footnotesize
\TODO{Fig.~4: GPU isolation. Planned content: latency distribution of a
concurrent CUDA workload run alone and run beside (i) the GPU baseline engine,
(ii) the fallback engine, and (iii) the strict DLA engine, plus the perception
latency distribution in each case. Requires the co-execution experiment of
Sec.~\ref{sec:rq4}; no data collected yet.}}}
\caption{Placeholder for the GPU-isolation result. The offloading decision of
this paper rests on this measurement together with Table~\ref{tab:energy}.}
\label{fig:isolation}
\end{figure}


RQ1--RQ3 establish the price: $1.72$ points of $\mathrm{mAP}_{50\text{--}95}$,
$7.36\times$ engine-boundary latency, and, in the current reference pipeline,
$2.94\times$ application latency. The paper's claim is that this price is
repaid in two currencies that matter on a power-capped robot with an
oversubscribed GPU, and this subsection is where that claim is tested. It is
not yet tested. We give the protocol, the instrumentation, and the table
skeleton, and mark every cell that requires data.

\paragraph*{Energy per frame}
Board power is logged after thermal stabilization in a fixed power mode, at
batch one, on the ordered image set, with the pipeline saturated and with it
running at a fixed frame rate; both total board energy per frame and
idle-subtracted energy per frame are retained, since the idle draw of the rest
of the module is not attributable to perception. The comparison of interest is
the strict engine against the GPU baseline at equal delivered frame rate, not at
equal utilization, because a robot's camera rate is fixed and the question is
what each placement costs per frame at that rate. Table~\ref{tab:energy} is the
intended form.

\begin{table}[t]
\caption{Energy accounting. Every cell requires the board-power measurement
described in Sec.~\ref{sec:rq4}; no energy data has been collected.}
\label{tab:energy}
\centering
\scriptsize
\setlength{\tabcolsep}{2.2pt}
\begin{tabular}{llccc}
\toprule
Power mode & Execution & Mean board & $E$/frame & $E$/frame \\
           &           & power (W)  & (mJ)      & idle-sub. (mJ) \\
\midrule
\TODO{mode A} & 11n / GPU       & \TODO{} & \TODO{} & \TODO{} \\
              & 11n / DLA+GPU   & \TODO{} & \TODO{} & \TODO{} \\
              & 11-DLA-n / DLA  & \TODO{} & \TODO{} & \TODO{} \\
\addlinespace
\TODO{mode B} & 11n / GPU       & \TODO{} & \TODO{} & \TODO{} \\
              & 11n / DLA+GPU   & \TODO{} & \TODO{} & \TODO{} \\
              & 11-DLA-n / DLA  & \TODO{} & \TODO{} & \TODO{} \\
\bottomrule
\end{tabular}
\end{table}

\paragraph*{GPU availability under contention}
Isolation is measured at two levels. First, a controlled CUDA workload runs
concurrently with perception, and its standalone latency distribution is
compared against its distribution beside the GPU baseline, beside the fallback
build, and beside the strict engine; this isolates contention and is what
Fig.~\ref{fig:isolation} will report. The strict engine is the only
configuration that should leave the concurrent workload undisturbed, and the
fallback build is the interesting control, since it occupies the GPU for only
$1.31\%$ of its engine time yet holds a stream, reformat kernels, and
$22.80$~MiB of context memory throughout. Second, the same engines are placed in
a ROS~2 perception node replaying a timestamped camera sequence at its
acquisition rate beside a fixed GPU-accelerated robotics workload
\cite{macenski2022ros2}; we report perception deadline misses, callback jitter,
and the slowdown of the concurrent node. The controlled test supports causal
attribution and the replay tests whether the difference is operationally
meaningful.
\TODO{all RQ4 results: concurrent-workload latency distributions, perception
deadline misses and jitter, and the slowdown of the co-running node, in each of
the three placements}

\subsection{Ablations and Validity}
\label{sec:validity}

\paragraph*{Ablations}
Two ablations explain the accuracy result of Sec.~\ref{sec:rq1} and one explains
the deployment result. \TODO{(i) the retrofit before and after the second
training phase, isolating how much of the retained accuracy comes from
weight-aware initialization versus fine-tuning; (ii) transferred versus random
initialization of the substituted block under an identical schedule; (iii) host
decoding on CPU versus GPU for the strict engine, which is the one configuration
where using the GPU for decoding partly defeats the purpose of the offload and
must therefore be quantified rather than assumed} An ablation that reintroduces
an inadmissible operator would be an accuracy control only and would not be
labeled accelerator-compatible.

\paragraph*{Threats to validity}
Every number in Table~\ref{tab:engine} comes from a single 30~s
\texttt{trtexec} launch with generated input buffers, and every number in
Table~\ref{tab:e2e} comes from one repeated image in one process launch;
neither samples content-dependent candidate counts, run-to-run variation, or
independent process launches, and the $6.4$~ms cross-process anomaly of
Sec.~\ref{sec:rq3} shows that the latter matters. Final timing estimates must
aggregate independent launches over an ordered image set with recorded power
mode, clocks, and thermal state. Accuracy is measured on the PyTorch path, so
Table~\ref{tab:accuracy} does not describe the deployed engine. Strict
compilation establishes \emph{placement}, not numerical equivalence: TensorRT
warns that its DLA softmax optimization is enabled for our channel softmax
because the volume of the non-axis dimensions exceeds one and that this may
introduce a small error in the output, so agreement must be measured on the
serialized engine rather than inferred from the build. The remaining
limitations are fixed input shapes, host-side post-processing, dependence on one
TensorRT version and one module, and a single detector scale; the formulation of
Sec.~\ref{sec:contract} is contract-agnostic but all evidence here comes from
one execution contract.
